%=====================================================================
\chapter{Functional Dependencies}\label{ch:fds}
%=====================================================================
\locator{5}{Part V --- Getting the Design Right}

\begin{outcomes}
By the end of this chapter you will be able to:
\begin{tight}
  \item \textbf{Define} a functional dependency and \textbf{read} one written
        as \fd{X}{Y}.
  \item \textbf{State} Armstrong's three axioms and \textbf{derive} new
        dependencies from given ones.
  \item \textbf{Compute} the closure $X^{+}$ of a set of attributes, step by
        step.
  \item \textbf{Use} attribute closure to test whether a set of attributes is
        a superkey or a candidate key.
  \item \textbf{Reduce} a set of dependencies to a minimal cover.
\end{tight}
\end{outcomes}

\begin{prereq}
\chref{relational} for keys. Appendix~B for sets and the meaning of
``function'' --- the Prerequisite Self-Check's question~12 is the whole
foundation of this chapter, and if it was uncomfortable, read Appendix~B
first, not during.
\end{prereq}

\begin{keyterms}
functional dependency $\bullet$ determinant $\bullet$ dependent $\bullet$
trivial dependency $\bullet$ Armstrong's axioms $\bullet$ reflexivity
$\bullet$ augmentation $\bullet$ transitivity $\bullet$ union rule $\bullet$
decomposition rule $\bullet$ attribute closure $\bullet$ $X^{+}$ $\bullet$
superkey $\bullet$ candidate key $\bullet$ minimal
cover $\bullet$ extraneous attribute $\bullet$ redundant dependency
\end{keyterms}

\section{One Table That Everyone Has Built}

Every university has a spreadsheet like this one. It is the first thing a
department builds when it is asked for ``all the results in one place'', and it
is the table this Part is about:

\medskip
\relation{report}{roll\_no, name, programme, course\_code, title,
credit\_hours, dept\_code, dept\_name, semester, section\_no, room, emp\_id,
instructor, grade}

\medskip
Fourteen columns, one row per enrolment, and on the seeded data 60{,}037 rows.
It answers every question in Part~IV without a single join, which is why
people build it. \chref{normalforms} measures what it costs. This chapter
supplies the vocabulary for saying \emph{why} it costs anything --- and the
vocabulary is one idea.

\begin{definitionbox}[Functional dependency]
A \term{functional dependency} \fd{X}{Y}, read ``$X$ determines $Y$'', holds in
a relation when any two tuples that agree on $X$ must also agree on $Y$.

\smallskip
$X$ is the \term{determinant}; $Y$ is the \term{dependent}. Both are sets of
attributes.

\smallskip
It is the word ``function'' in its mathematical sense, from Appendix~B: give
me the value of $X$ and there is exactly one value of $Y$ it can be paired
with. \fd{roll\_no}{name} says a roll number has exactly one name.
\end{definitionbox}

\begin{examplebox}[The dependencies in \textit{report}, found by asking]
As with keys in \chref{ermodel}, these come from the university, not from the
data. Each one is a rule somebody can confirm:

\medskip
\begin{tabular}{@{}L{6.6cm}L{7.4cm}@{}}
\textbf{Dependency} & \textbf{Because} \\
\fd{roll\_no}{name, programme} & a student has one name and is in one
  programme \\
\fd{course\_code}{title, credit\_hours, dept\_code} & a course is defined
  once \\
\fd{dept\_code}{dept\_name} & a department has one name \\
\fd{course\_code, semester, section\_no}{room, emp\_id} & a section meets in
  one room with one instructor \\
\fd{emp\_id}{instructor} & an employee has one name \\
\fd{roll\_no, course\_code, semester, section\_no}{grade} & one grade per
  enrolment \\
\end{tabular}

\smallskip
Six rules. Notice that each one is a sentence the examinations officer would
agree with immediately --- and that each one says the dependent is stored in
\textit{report} \emph{once per enrolment}, when it is one fact. Ayesha's name
is in the table once for every section she has ever taken.
\end{examplebox}

\begin{alertbox}[An instance can disprove a dependency, never prove one]
\chref{relational} said this about keys, and it is the same fact, because a
key is a dependency (below). On the seeded data, \textit{name} happens to
determine \textit{programme} for some names. That is coincidence, not a rule
--- two students called Hira Nadeem can be in different programmes, and the
next admission may be the one that proves it.

\smallskip
A query \emph{can} disprove a dependency, and it is worth knowing:

\begin{lstlisting}[language=sql]
-- Does course_code -> title hold in this data?
SELECT course_code FROM report
 GROUP BY course_code HAVING COUNT(DISTINCT title) > 1;
-- Any row returned is a counter-example. None returned proves nothing.
\end{lstlisting}
\end{alertbox}

\section{Deriving New Dependencies}

\begin{definitionbox}[Armstrong's axioms]
For sets of attributes $X$, $Y$, $Z$:

\smallskip
\textbf{Reflexivity.} If $Y \subseteq X$ then \fd{X}{Y}. (A set determines its
own parts. Such a dependency is \term{trivial}.)

\smallskip
\textbf{Augmentation.} If \fd{X}{Y} then \fd{XZ}{YZ}. (Adding the same
attributes to both sides keeps it true.)

\smallskip
\textbf{Transitivity.} If \fd{X}{Y} and \fd{Y}{Z} then \fd{X}{Z}.

\smallskip
These three are \emph{sound} (they never derive anything false) and
\emph{complete} (they derive everything that follows). Two derived rules save
time: \textbf{union} --- \fd{X}{Y} and \fd{X}{Z} give \fd{X}{YZ}; and
\textbf{decomposition} --- \fd{X}{YZ} gives \fd{X}{Y} and \fd{X}{Z}.
\end{definitionbox}

\begin{examplebox}[One transitive dependency, and why it is the important one]
From the table:
\begin{tight}
  \item \fd{course\_code}{dept\_code} \hfill (given, by decomposition)
  \item \fd{dept\_code}{dept\_name} \hfill (given)
  \item therefore \fd{course\_code}{dept\_name} \hfill (transitivity)
\end{tight}

\smallskip
Nobody wrote that third rule down, and it is true. It is also the one that
causes trouble: \textit{dept\_name} depends on \textit{course\_code} only
\emph{through} \textit{dept\_code}. Renaming the department means changing
every row of every course it offers --- \chref{normalforms} calls that a
transitive dependency and third normal form forbids it.
\end{examplebox}

\section{Attribute Closure}

Deriving dependencies one axiom at a time is how proofs are written. The
practical question is narrower: \emph{given $X$, what does it determine?} There
is an algorithm, and it is the most useful thing in Part~V.

\begin{definitionbox}[Attribute closure $X^{+}$]
The \term{closure} $X^{+}$ of a set of attributes $X$ under a set of
dependencies $F$ is the set of every attribute that $X$ determines.

\smallskip
\textbf{Algorithm.}
\begin{tightnum}
  \item Start with $X^{+} = X$.
  \item Repeat: for each dependency \fd{A}{B} in $F$, if $A \subseteq X^{+}$,
        add $B$ to $X^{+}$.
  \item Stop when a full pass adds nothing.
\end{tightnum}

\smallskip
It always terminates --- there are finitely many attributes and $X^{+}$ only
grows.
\end{definitionbox}

\begin{worked}[$\{\mathit{course\_code}\}^{+}$, then the key of \textit{report}]
Using the six dependencies from the examplebox, abbreviated
$F$.

\medskip
\textbf{Part A --- what does \textit{course\_code} determine?}

\smallskip
\begin{tabular}{@{}L{1.4cm}L{6.0cm}L{6.6cm}@{}}
\textbf{pass} & \textbf{dependency used} & \textbf{$X^{+}$ is now} \\
start & --- & \{course\_code\} \\
1 & \fd{course\_code}{title, credit\_hours, dept\_code} & \{course\_code,
  title, credit\_hours, dept\_code\} \\
1 & \fd{dept\_code}{dept\_name} --- its left side is now inside & \dots\ plus
  dept\_name \\
2 & nothing new applies & \textbf{stop} \\
\end{tabular}

\smallskip
$\{\mathit{course\_code}\}^{+} = \{$course\_code, title, credit\_hours,
dept\_code, dept\_name$\}$. Five of the fourteen. \textit{course\_code} is not a
key of \textit{report}.

\medskip
\textbf{Part B --- the key.} Try $K = \{$roll\_no, course\_code, semester,
section\_no$\}$:
\begin{tight}
  \item \fd{roll\_no}{name, programme} adds 2;
  \item \fd{course\_code}{title, credit\_hours, dept\_code} adds 3, and then
        \fd{dept\_code}{dept\_name} adds 1;
  \item \fd{course\_code, semester, section\_no}{room, emp\_id} adds 2, and
        then \fd{emp\_id}{instructor} adds 1;
  \item \fd{roll\_no, course\_code, semester, section\_no}{grade} adds 1.
\end{tight}
$4 + 2 + 3 + 1 + 2 + 1 + 1 = 14$. $K^{+}$ is every attribute: \textbf{$K$ is a
superkey}.

\medskip
\textbf{Part C --- is it a candidate key?} Remove each attribute in turn and
recompute. Without \textit{roll\_no}, nothing determines \textit{name}; without
\textit{course\_code}, nothing determines \textit{title}; without
\textit{semester} or \textit{section\_no}, nothing determines \textit{room}. No
proper subset is a superkey, so \textbf{$K$ is a candidate key} --- and it is
exactly \chref{mapping}'s primary key of \textit{enrolment}, found here from
the dependencies alone.
\end{worked}

\dsfig{%
\begin{tikzpicture}[
  atr/.style={rounded corners=2pt, draw=#1, line width=0.8pt, fill=#1!10,
             text=#1!60!black, font=\scriptsize\ttfamily, inner sep=2.5pt,
             minimum height=0.48cm},
  ps/.style={font=\scriptsize\bfseries, text=primary, anchor=east},
  rule/.style={font=\scriptsize\itshape, text=neutral, anchor=west},
  arr/.style={-{Stealth[length=2mm]}, primary!45, line width=0.8pt}
]
  % start
  \node[ps] at (-0.25,1.8) {start};
  \node[atr=primary, anchor=west] (s0) at (0,1.8) {course\_code};

  % pass 1, first rule
  \node[ps] at (-0.25,0.9) {pass 1};
  \node[atr=primary, anchor=west] (a1) at (0,0.9) {course\_code};
  \node[atr=secondary, anchor=west] (a2) at (a1.east) {title};
  \node[atr=secondary, anchor=west] (a3) at (a2.east) {credit\_hours};
  \node[atr=secondary, anchor=west] (a4) at (a3.east) {dept\_code};
  \node[rule] at (7.6,0.9) {\fd{course\_code}{title, credit\_hours,
       dept\_code}};

  % pass 1, second rule, now applicable
  \node[ps] at (-0.25,0.0) {pass 1};
  \node[atr=primary, anchor=west] (b1) at (0,0.0) {course\_code};
  \node[atr=secondary, anchor=west] (b2) at (b1.east) {title};
  \node[atr=secondary, anchor=west] (b3) at (b2.east) {credit\_hours};
  \node[atr=secondary, anchor=west] (b4) at (b3.east) {dept\_code};
  \node[atr=greenhl, anchor=west] (b5) at (b4.east) {dept\_name};
  \node[rule] at (7.6,0.0) {\fd{dept\_code}{dept\_name} --- its left side
       is now inside};

  % pass 2
  \node[ps] at (-0.25,-0.9) {pass 2};
  \node[font=\scriptsize\itshape, text=accent, anchor=west] at (0,-0.9)
       {no dependency adds anything new $\Rightarrow$ \textbf{stop}};

  \draw[arr] (s0.south) ++(0,-0.02) -- ([yshift=1pt]a1.north);
  \draw[arr] (a1.south) ++(0,-0.02) -- ([yshift=1pt]b1.north);
\end{tikzpicture}}%
{$\{\mathit{course\_code}\}^{+}$, one step per row. Each step applies a
dependency whose left side is already inside the set; the second only becomes
applicable because the first added \textit{dept\_code}. The closure stops at
five attributes --- fewer than fourteen, so \textit{course\_code} is not a key
of \textit{report}.}{closure}

\begin{conceptbox}[Why closure is the tool for everything in Part V]
Three questions, one algorithm:

\begin{tight}
  \item \textbf{Is $X$ a superkey?} Compute $X^{+}$; it is if $X^{+}$ is every
        attribute.
  \item \textbf{Does \fd{X}{Y} follow from $F$?} Compute $X^{+}$; it does if
        $Y \subseteq X^{+}$. No axioms needed.
  \item \textbf{Is an attribute in a dependency redundant?} Remove it and
        recompute; if the closure is unchanged, it was. That is the minimal
        cover, next.
\end{tight}

\smallskip
Every normal form in \chref{normalforms} is a test you run with closures, and
\chref{decomposition}'s lossless-join test is one too.
\end{conceptbox}

\section{Minimal Cover}

\begin{definitionbox}[Minimal cover]
A \term{minimal cover} of $F$ is an equivalent set of dependencies --- one with
the same closure --- in which:
\begin{tightnum}
  \item every right-hand side is a single attribute;
  \item no left-hand side has an \term{extraneous attribute} (one that could be
        removed without changing what the dependency implies);
  \item no dependency is \term{redundant} (derivable from the others).
\end{tightnum}

\smallskip
It is the smallest statement of the same rules. \chref{normalforms}'s 3NF
synthesis algorithm starts from one.
\end{definitionbox}

\begin{worked}[A minimal cover, in three steps]
$F = \{$ \fd{A}{BC}, \fd{B}{C}, \fd{AB}{D} $\}$.

\medskip
\textbf{Step 1 --- single right-hand sides.} Decompose:
\fd{A}{B}, \fd{A}{C}, \fd{B}{C}, \fd{AB}{D}.

\medskip
\textbf{Step 2 --- extraneous left-hand attributes.} Only \fd{AB}{D} has two.
Is $B$ extraneous? Compute $\{A\}^{+}$: $A \rightarrow B$ gives $B$, so
$\{A\}^{+} \ni B$, and then \fd{AB}{D} applies. $\{A\}^{+} = \{A,B,C,D\}$,
which contains $D$. So \fd{A}{D} already holds: \textbf{$B$ is extraneous},
and \fd{AB}{D} becomes \fd{A}{D}.

\medskip
\textbf{Step 3 --- redundant dependencies.} Is \fd{A}{C} redundant? Remove it
and compute $\{A\}^{+}$ from the rest: $A \to B$, $B \to C$. It contains $C$.
\textbf{Redundant} --- it follows by transitivity. Remove it.

\medskip
\textbf{Result.} $\{$ \fd{A}{B}, \fd{B}{C}, \fd{A}{D} $\}$. Three dependencies
instead of four, each as small as it can be, and the same closure as the
original. Check: $\{A\}^{+} = \{A,B,C,D\}$ in both.
\end{worked}

\begin{pitfall}
\begin{tight}
  \item \textbf{Reading a dependency off sample data.} Data disproves; it never
        proves. Ask.
  \item \textbf{Reading \fd{X}{Y} backwards.} \fd{roll\_no}{name} does not
        give \fd{name}{roll\_no}; two students share a name.
  \item \textbf{Stopping closure after one pass.} A dependency whose left side
        becomes available in pass 1 must still be applied. Repeat until a pass
        adds nothing.
  \item \textbf{Calling a superkey a candidate key.} Test every proper subset.
  \item \textbf{Splitting a left-hand side.} \fd{AB}{C} does \emph{not} give
        \fd{A}{C}. Only right-hand sides decompose.
  \item \textbf{Forgetting the transitive dependencies nobody wrote down.}
        Closure finds them; that is what it is for.
\end{tight}
\end{pitfall}

\begin{lab}[Build the report table, then find its dependencies and its key]
\begin{lstlisting}[language=bash]
# --- 1. The seeded database ----------------------------------------
python code/schema/seed.py university.db
sqlite3 university.db
\end{lstlisting}
\begin{lstlisting}[language=sql]
-- --- 2. Build the table everybody builds ------------------------
CREATE TABLE report AS
SELECT s.roll_no, s.name, s.programme,
       c.course_code, c.title, c.credit_hours,
       d.dept_code, d.name AS dept_name,
       e.semester, e.section_no, x.room, x.emp_id,
       i.name AS instructor, e.grade
  FROM enrolment e
  JOIN student    s ON s.roll_no     = e.roll_no
  JOIN course     c ON c.course_code = e.course_code
  JOIN department d ON d.dept_code   = c.dept_code
  JOIN section    x ON x.course_code = e.course_code
                   AND x.semester    = e.semester
                   AND x.section_no  = e.section_no
  LEFT JOIN instructor i ON i.emp_id = x.emp_id;
SELECT COUNT(*) FROM report;

-- --- 3. Test a dependency that should hold ----------------------
SELECT course_code FROM report
 GROUP BY course_code HAVING COUNT(DISTINCT title) > 1;
-- No rows: no counter-example. Does that PROVE course_code -> title?

-- --- 4. Test one that should not --------------------------------
SELECT name FROM report
 GROUP BY name HAVING COUNT(DISTINCT roll_no) > 1 LIMIT 5;
-- Counter-examples: name does not determine roll_no. Settled.

-- --- 5. Count the copies of one fact ----------------------------
SELECT roll_no, name, COUNT(*) AS copies FROM report
 GROUP BY roll_no ORDER BY copies DESC LIMIT 3;
-- One student's name, stored this many times. Why?

-- --- 6. The transitive one -------------------------------------
SELECT dept_code, COUNT(*) AS rows_holding_its_name
  FROM report GROUP BY dept_code;
-- To rename one department, how many rows must change?
\end{lstlisting}
\begin{lstlisting}[language=bash]
# --- 7. On paper -----------------------------------------------------
# Compute {course_code}+ , {roll_no}+ and {emp_id}+ by the algorithm,
# one pass per line, from the six dependencies in this chapter.
# Then prove {roll_no, course_code, semester, section_no} is a
# candidate key: superkey first, then every proper subset.
\end{lstlisting}

\smallskip
\textbf{What to notice.} Steps 5 and 6 are the whole motivation for Part~V. One
fact --- a name, a department name --- is stored once for every enrolment it
touches, because the dependency that determines it has a smaller left side than
the table's key. Keep \textit{report}: \chref{normalforms} measures it against
the normalised schema, and \chref{decomposition} measures it again.
\end{lab}

\begin{checkpoint}
\begin{tightnum}
  \item Compute $\{\mathit{emp\_id}\}^{+}$ and
        $\{\mathit{course\_code}, \mathit{semester}, \mathit{section\_no}\}^{+}$
        under the six dependencies of \textit{report}, showing each pass.
  \item Using closure only --- no axioms --- decide whether
        \fd{roll\_no, course\_code}{dept\_name} follows from the six
        dependencies, and show the closure that settles it.
  \item Find a minimal cover of
        $\{$ \fd{A}{B}, \fd{A}{C}, \fd{B}{C}, \fd{AC}{D} $\}$, showing all
        three steps and checking that the result has the same closure of $A$.
\end{tightnum}
\end{checkpoint}

\begin{chaptersummary}
\begin{tight}
  \item \fd{X}{Y} holds when tuples agreeing on $X$ must agree on $Y$. It is a
        function, in Appendix~B's sense, from $X$-values to $Y$-values.
  \item Dependencies come from the university. An instance can disprove one
        with a \sqlkw{GROUP BY} \dots\ \sqlkw{HAVING} \sqlkw{COUNT(DISTINCT)}
        query; it can never prove one.
  \item \textbf{Armstrong's axioms} --- reflexivity, augmentation,
        transitivity --- are sound and complete. Union and decomposition
        follow from them.
  \item Only right-hand sides decompose. \fd{AB}{C} does not give
        \fd{A}{C}.
  \item The \textbf{closure} $X^{+}$ is everything $X$ determines: start with
        $X$, apply every dependency whose left side is inside, repeat until a
        pass adds nothing.
  \item Closure answers three questions: is $X$ a superkey, does a dependency
        follow, is an attribute redundant.
  \item The key of \textit{report} is \{roll\_no, course\_code, semester,
        section\_no\} --- the primary key of \textit{enrolment}, found from
        the dependencies alone.
  \item A \textbf{minimal cover} has single-attribute right sides, no
        extraneous left-side attributes, and no redundant dependencies.
  \item In \textit{report}, every fact whose determinant is smaller than the
        key is stored once per enrolment. That is what \chref{normalforms}
        removes, and measures.
\end{tight}
\end{chaptersummary}

\begin{reviewq}
  \item \textbf{(MCQ)} \fd{AB}{C} and \fd{C}{D} together imply:
        (a)~\fd{A}{D}; (b)~\fd{AB}{D}; (c)~\fd{D}{AB}; (d)~\fd{C}{AB}.
  \item \textbf{(MCQ)} A dependency \fd{X}{Y} with $Y \subseteq X$ is:
        (a)~transitive; (b)~trivial; (c)~redundant; (d)~a key.
  \item \textbf{(Short)} Why can a query find a counter-example to a
        dependency but never prove one holds?
  \item \textbf{(Short)} State the closure algorithm in three lines, and say
        why it must terminate.
  \item \textbf{(Short)} Show by a two-row example that \fd{AB}{C} does not
        imply \fd{A}{C}.
  \item \textbf{(Applied)} A relation \relation{booking}{room, date, slot,
        course\_code, emp\_id, capacity} has the rules: a room has a fixed
        capacity; a room, date and slot is booked for one course; a course has
        one instructor per booking. Write the dependencies, compute the closure
        of \{room, date, slot\}, and find every candidate key.
  \item \textbf{(Applied)} Using the six dependencies of \textit{report}, list
        every attribute whose determinant is a proper subset of the key, and for
        each one compute from the seeded data (step~5 of the laboratory) how
        many copies of a single fact it causes. Which one costs the most rows?
\end{reviewq}
